\begin{abstract}
Báo cáo này trình bày thiết kế, hiện thực và đánh giá thực nghiệm một đường ống dữ liệu tài chính tự động (Financial Data Engineering Pipeline) thu thập, chuẩn hóa và kiểm soát chất lượng các bản thuyết minh yếu tố rủi ro (\textit{Item 1A. Risk Factors}) từ báo cáo thường niên Form 10-K của Ủy ban Chứng khoán và Giao dịch Hoa Kỳ (SEC EDGAR). Hệ thống áp dụng kiến trúc hồ chứa dữ liệu 4 tầng Medallion (Bronze, Silver, Gold và Curated) hoàn toàn trên hệ thống tệp tin cục bộ định dạng Apache Parquet nén Snappy, loại bỏ sự phụ thuộc vào các cụm hạ tầng phân tán cồng kềnh.

Dự án xác lập ba đóng góp kỹ thuật trọng tâm: thứ nhất, \textbf{cơ chế cảm biến chênh lệch (Delta Sensor)} theo dõi mã kiểm soát \texttt{accessionNumber} từ API siêu dữ liệu SEC để phát hiện trạng thái cập nhật và ngắt ngắn chuỗi tác vụ dư thừa trong 3,09 giây với 0 byte băng thông lãng phí; thứ hai, \textbf{quy trình bóc tách lai DOM--Regex} loại bỏ triệt để rác siêu văn bản iXBRL và các thẻ kịch bản lồng nhau, vượt bẫy bảng mục lục (\textit{Table of Contents Trap}) qua ngưỡng 5.000 ký tự và chuẩn hóa mã định danh pháp lý CIK 10 chữ số chống lỗi ép kiểu số thực; thứ ba, \textbf{khắc phục lỗi suy giảm chất lượng dữ liệu âm thầm (\textit{Truncation Defect})} phát hiện 98,1\% phân đoạn đếm từ bị cắt cụt ngữ cảnh và tái thiết kế bằng giải thuật phân đoạn nhận biết bộ tách từ (\textit{tokenizer-aware chunking}) $\le$240 tokens, được khóa chặt bởi hợp đồng dữ liệu tự động trong bộ 24 ca kiểm thử Pytest.

Kết quả thực nghiệm xác nhận hệ thống rút giảm 94,3\% dung lượng lưu trữ (từ 22,51 MB dữ liệu thô xuống 1,28 MB dữ liệu sạch), đảm bảo tính tất định và khả năng tái lập 100\%, đồng thời phục hồi 50\% độ chính xác truy xuất thông tin (Hit@15) trong các vùng điểm mù ngữ cảnh phục vụ các mô hình phân tích hạ nguồn.
\end{abstract}
